\chapter{Navigation, zoom et comparaison}
\label{chap:navigation}
\section{Passer d’une image à l’autre}
Cliquer sur une vignette pour ouvrir sa photo. Le contour de la vignette et
son nom en gras identifient l’image effectivement ouverte, indépendamment des
cases de sélection multiple. Si le ruban est replié, les boutons précédent
et suivant apparaissent en bas à droite de la zone image. Leur survol indique
le nom de la cible ; ils ne dépassent pas les extrémités du dossier.

Le bouton centré au-dessus du ruban le replie. Lorsqu’il est masqué, ce bouton
revient en bas de la vue et permet de le rouvrir. Son état est mémorisé.
Les flèches du clavier ne sont pas des raccourcis globaux de changement de photo.

\section{Examiner les détails}
Cliquer dans la photo pour passer du mode ajusté à un agrandissement, puis
cliquer de nouveau sans déplacer la souris pour revenir au mode ajusté.
Quand l’image dépasse la vue, maintenir le bouton gauche et déplacer la souris
pour parcourir les détails. Un clic dans une marge vide ne déclenche pas le zoom.

\paragraph{Clavier et molette.}
\touche{Ctrl + +} agrandit ; \touche{Ctrl + $-$} réduit.
\touche{Ctrl + 0} ajuste à la vue et \touche{Ctrl + 1} affiche à 100\,\%.
\touche{Ctrl + molette} zoome autour de la souris. Sans Ctrl, la molette déplace
l’image agrandie ; \touche{Maj + molette verticale} permet un déplacement horizontal.
Les boutons correspondants restent disponibles au-dessus de la photo.

\attention{100\,\% signifie un pixel de l’aperçu par unité de la vue ;
l’aperçu interactif peut être réduit par rapport à la source pleine résolution.
L’export est reconstruit indépendamment à sa résolution prévue.}

\section{Afficher un comparatif partagé}
Cliquer sur \emph{Avant/après} ou utiliser \touche{Ctrl + Maj + C} lorsque
les deux aperçus sont prêts. L’origine est à gauche et l’état courant à droite.
Déplacer le séparateur pour choisir la largeur des deux zones. Les deux côtés
partagent le même zoom et le même déplacement.

\section{Voir temporairement l’original}
Donner le focus à la vue image, puis maintenir \touche{Espace}. Relâcher revient
à la présentation précédente, y compris au comparatif partagé s’il était actif.
Le bouton \emph{Origine} remplit le même rôle tant que le clic est maintenu.
La commande ne s’applique pas à une saisie dans un champ numérique.
Un changement de photo réinitialise la vue et la comparaison.

\attention{L’origine est la référence de développement exposée par IPR,
pas la mosaïque RAW brute ni nécessairement le JPEG du boîtier. Modifier
le profil de développement peut modifier cette référence.}

\section{Menu contextuel de la photo}
Un clic droit dans la vue ouvre les actions de la photo : export et
réinitialisation des retouches, ainsi que les commandes de fichier.
Le clic droit ne déclenche pas le zoom. Avec le focus dans la vue,
\touche{Maj + F10} ou la touche Menu ouvre aussi ce contexte.
La réinitialisation demande confirmation ; elle n’a pas d’annulation générale.
Les actions Enregistrer et Dupliquer restent désactivées tant que leur
raccordement moteur manque.
